\section{Chapitre~\ref{sec:chemoutput}. La sortie chimique}

Les deux fils vers le cœur et leurs vitesses différentes, la contre-réaction dans les cascades hormonales, le code par fréquence des bouffées, le générateur circadien et son recalage par la lumière sont mesurés directement; la limite $K<8$ est un théorème de l'automatique démontré dans le chapitre; les pulsations engendrées par la rétroaction de la cascade elle-même sont une hypothèse soutenue par une expérience solide.

{\footnotesize
\setlength{\tabcolsep}{3pt}\renewcommand{\arraystretch}{1.15}
\begin{longtable}{>{\raggedright}p{0.5cm}>{\raggedright}p{4.0cm}>{\raggedright}p{5.6cm}>{\raggedright}p{2.3cm}>{\raggedright\arraybackslash}p{1.7cm}}
\toprule
N\textsuperscript{o} & Affirmation & Expérience et ce qui est mesuré & Source & Statut \\
\midrule\endfirsthead
\toprule
N\textsuperscript{o} & Affirmation & Expérience et ce qui est mesuré & Source & Statut \\
\midrule\endhead
1 & Le stimulateur du cœur bat de lui-même environ $100$ fois par minute; au repos, le frein est enfoncé, et la fréquence est en général de l'ordre de $60$--$70$ & Humains, blocage complet des deux fils vers le cœur: la fréquence propre dépasse la fréquence de repos et décroît avec l'âge; chez l'adulte, elle est de l'ordre de $100$ battements par minute. Au repos, c'est donc le frein qui domine. La fréquence de repos de $60$--$70$ est une valeur typique, non citée séparément. & \cite{JoseCollison1970ihr} & mesuré \\
2 & L'accélérateur \nt{NORA} et le frein \nt{ACET} sont des filtres passe-bas, le frein est plus rapide~\eqref{eq:gasbrake} & Chien, stimulation modulée en fréquence du nerf vague ou du nerf sympathique cardiaque, et fonction de transfert vers la fréquence auriculaire: les deux voies sont des filtres passe-bas; la voie sympathique a une fréquence de coupure bien plus basse et un retard pur d'environ $1.7\,\text{s}$ en plus. Les caractéristiques dépendent du tonus moyen, si bien que la formule linéaire~\eqref{eq:gasbrake} est une approximation. & \cite{Berger1989} & mesuré \\
3 & Le frein est rapide parce qu'\nt{ACET} ouvre le canal potassique sans second messager, alors que \nt{NORA} passe par une chaîne de réactions & Fragments de membrane de cellules auriculaires: le canal potassique ouvert par l'acétylcholine est ouvert par les sous-unités $\beta\gamma$ de la protéine G couplée au récepteur, sans second messager intracellulaire. La voie de \nt{NORA} par l'AMPc n'est pas citée ici. & \cite{Logothetis1987gbg} & mesuré \\
4 & Le sang fait le tour du corps en une minute environ; \nt{ADRE} circulant atteint tous les organes munis du récepteur & Estimation générale d'ordre de grandeur; formulée ainsi dans le chapitre, sans source citée. & --- & estimation \\
5 & La cascade hormonale est bouclée par une contre-réaction: l'hormone finale freine les premiers étages & Revue d'expériences: les hormones corticosurrénaliennes répriment la libération d'ACTH par l'hypophyse et de l'hormone de libération par l'hypothalamus dans trois gammes de temps: rapide, intermédiaire et lente. & \cite{KellerWood1984fb} & mesuré \\
6 & Trois étages identiques sont stables pour $K<8$~\eqref{eq:cascadestable}; à la limite, la période vaut $2\pi\tau/\sqrt3\approx3.6\tau$ & Démontré dans le chapitre: à la pulsation $\sqrt3/\tau$, trois étages donnent un déphasage de $180^\circ$ et une atténuation d'un facteur $8$. & démonstration dans le chapitre & théorie \\
7 & Erreur sur le niveau $1/(1+K)$: un tel régulateur ne tient donc pas mieux qu'environ $10\,\%$ (proposition~\ref{prop:cascade}) & Conséquence de la ligne~6 pour une rétroaction proportionnelle. L'axe réel a une rétroaction sur plusieurs étages et dans plusieurs gammes de temps (ligne~5): la limite se rapporte donc au modèle~\eqref{eq:cascade} et n'est pas mesurée. & démonstration dans le chapitre & théorie \\
8 & Pour $K=4$ et $7$, le niveau se stabilise; pour $K=12$, pulsations régulières de période $3.7$--$3.9\,\tau$ (fig.~\ref{fig:cascade}) & Calcul avec \texttt{models/\allowbreak chem\_\allowbreak models.py}: pour $K=12$, périodes successives $3.9$, $3.7$, $3.8$, $3.7\,\tau$; pour $K=4$, amplitude en fin de calcul $0.003$. & \texttt{models/\allowbreak chem\_\allowbreak models.py} & modèle \\
9 & L'hormone du stress sort par pulsations environ une fois par heure; les pulsations naissent de la rétroaction elle-même, sans générateur séparé & Rats, perfusion constante de l'hormone de libération: l'ACTH et la corticostérone oscillent à une fréquence physiologique, proche d'une heure; une entrée rythmique venant de l'hypothalamus n'est pas nécessaire aux pulsations. Un modèle hypophyse--surrénale à rétroaction retardée produit de telles oscillations. & \cite{Walker2012glc, Walker2010} & hypothèse \\
10 & Le destinataire lit la fréquence des bouffées, non le niveau & Singes privés de leur propre libération de l'hormone de libération des gonadotrophines: une perfusion constante ne rétablit pas la sécrétion des hormones hypophysaires, des bouffées toutes les heures la rétablissent; le retour à la perfusion constante éteint de nouveau la réponse. La fatigue du récepteur vue comme filtre passe-haut est une interprétation. & \cite{Belchetz1978} & mesuré \\
11 & Des fréquences de bouffées différentes agissent différemment sur les diverses cellules d'une même glande & Mêmes singes: porter les bouffées à $2$--$5$ par heure fait baisser les deux hormones hypophysaires; les espacer à une toutes les $3$~heures fait baisser l'une (LH) et monter l'autre (FSH), si bien que leur rapport est fixé par la fréquence. & \cite{Wildt1981gnrh} & mesuré \\
12 & Le rythme circadien naît d'une contre-réaction intracellulaire retardée de plusieurs heures & Revue: les protéines des gènes d'horloge répriment avec retard la transcription de leurs propres gènes; la boucle transcription--traduction donne une période d'environ un jour. & \cite{TakahashiJS2017} & mesuré \\
13 & Le générateur principal est un groupe de quelques dizaines de milliers de neurones qui synchronise le reste du corps & Revue: le noyau suprachiasmatique est le générateur circadien principal; ses neurones isolés en culture sont des générateurs autonomes, et le réseau les synchronise; chez la souris, environ $10^4$ neurones de chaque côté. & \cite{Welsh2010scn} & mesuré \\
14 & La période propre chez l'humain est de $24.18$~heures & Humains en lumière faible, avec un horaire découplé du jour: période des rythmes de la mélatonine, de la température et du cortisol en moyenne de $24.18$~heures chez les jeunes comme chez les âgés, avec une faible dispersion. & \cite{Czeisler1999} & mesuré \\
15 & La lumière recale le générateur comme une boucle à verrouillage de phase~\eqref{eq:pll}; il faut un décalage d'environ $11$~min par jour & Humains, une seule exposition à une lumière vive de $6.7$~heures à différentes heures: courbe de réponse de phase de type~1, d'amplitude $5$~heures; retard avant le minimum de la température corporelle, avance après. L'équation~\eqref{eq:pll} est le modèle standard; $11$~min $=0.18$~heure est de l'arithmétique. & \cite{Khalsa2003prc} & mesuré \\
16 & Sans lumière, pas de verrouillage: le rythme dérive vers sa période propre & Personnes totalement aveugles, sans perception de la lumière, vivant en société normale: chez environ la moitié d'entre elles, les rythmes de la mélatonine et du cortisol suivent leur propre période, différente de celle du jour. & \cite{Sack1992blind} & mesuré \\
\bottomrule
\end{longtable}}
